\section{Invariants}\label{sec:invariants}

This section derives the invariants of one compressive branch from the slice equation (\ref{eq:slice}), for smooth Alfv\'enic fields and with no collisions, dissipation or sources acting on $G$; the value $\Lam=\infty$ is included. We write $\langle\cdot\rangle$ for the average over the periodic box, of volume $V$.

\subsection{The quadratic family}\label{inv:family}

Because the slices diagonalise $L$, a function $\varphi$ of $L$ acts on them as multiplication by $\varphi(u)$: the slice amplitude of $\varphi(L)G$ is $\varphi(u)\,a_u$. By (\ref{sl:parseval}), the quantity
\begin{equation}\label{eq:Qphi}
  \Qphi\equiv\int\dd^3r\,\fe{G}{\varphi(L)G}=\int\dd u\;\varphi(u)\int\dd^3r\,a_u^2=V\!\int\varphi(u)\,\mathcal E(u)\,\dd u,
  \qquad \mathcal E(u)\equiv\langle a_u^2\rangle ,
\end{equation}
weights the free energy of each slice with $\varphi(u)$; for unbounded $\varphi$ the last two expressions define $\Qphi$ whenever the integral converges.\footnote{For $\varphi(u)=u^n$ the integral converges when $\int\dd^3r\int(1+u^{2k})\,G^2F_M^{-1}\,\dd u$ is finite for an integer $k\ge n/2$. Then $|u|^n\le1+u^{2k}$, and $\int u^{2k}a_u^2\,\dd u=\fe{L^kG}{L^kG}$ is finite, because $L=UP$ (\ref{eq:LUP}) differs from multiplication by $u$ only through $P$, which changes $G$ by a multiple of $F_M$, a function with finite moments of every order. For $\varphi(u)=1/u$ the integral is a principal value, which exists when $G(\boldsymbol r,\cdot)$ is H\"older continuous near $u=0$, uniformly in $\boldsymbol r$, because $a_u$, by (\ref{eq:aunit}) a combination of $G$ and its Hilbert transform with coefficients smooth in $u$, is then H\"older continuous there too.} The slice energy $\mathcal E(u)$ is the distribution of free energy over phase velocity: $\tfrac12n_iT_i\,\mathcal E(u)\,\dd u$ is the free energy per unit volume in the slices with phase velocities in $[u,u+\dd u]$.

Each slice is a passive scalar carried by the divergence-free velocity $\boldsymbol U_u$, so its energy density obeys the conservation law $\partial_ta_u^2+\nabla\cdot(\boldsymbol U_u\,a_u^2)=0$. Integrating it over the periodic box, we find that the energy of every slice is conserved:
\begin{equation}\label{inv:energy}
  \frac{\dd\mathcal E(u)}{\dd t}=0 .
\end{equation}
The nonlinearity therefore moves no free energy between phase velocities, and $\Qphi$ is conserved for every weight $\varphi$ for which the integral in (\ref{eq:Qphi}) exists at $t=0$, either absolutely or, for $\varphi=1/u$, as a principal value. This is an infinite family of quadratic invariants.

Per-velocity quadratic invariants are known in the simpler settings of section~\ref{sl:remlit} \citep{PlunkParker2014,Plunk2015}, and, unstated, in the transformed equation of \citet[eq.~85]{Heninger2018}. \citet[\S5.1]{Kunz2015} remarked that the free energy of their collisionless pressure-anisotropic generalisation of KRMHD is one of infinitely many invariants, without specifying the others.

\subsection{The invariants \texorpdfstring{$W$, $\Gamma$ and $\Hph$}{W, Gamma and H}}\label{inv:cmm}

The free energy $W^\pm$ of each branch (CMM's B13) is conserved (S09), and CMM found two further quadratic invariants for each branch, the phase-space helicity $\Hph^\pm$ (their B18) and $\Gamma^\pm$ (their B25). Written in $u=v_\parallel/\vth$, with $G^\pm(v_\parallel)\,\dd v_\parallel=G^\pm(u)\,\dd u$, they are the members $\varphi=1$, $u$ and $1/u$ of (\ref{eq:Qphi}):
\begin{equation}\label{eq:WGH}
\begin{aligned}
  W^\pm&=\frac{n_iT_i}{2}\,\big\langle\fe{G^\pm}{G^\pm}\big\rangle=\frac{n_iT_i}{2}\int\mathcal E^\pm(u)\,\dd u,\\
  \Gamma^\pm&=\frac{\vth}{2}\,\big\langle\fe{G^\pm}{LG^\pm}\big\rangle=\frac{\vth}{2}\int u\,\mathcal E^\pm(u)\,\dd u,\\
  \Hph^\pm&=\frac{n_iT_i}{2\vth}\,\Big\langle\PVint\frac{(G^\pm)^2}{uF_M}\,\dd u\Big\rangle=\frac{n_iT_i}{2\vth}\,\PVint\frac{\mathcal E^\pm(u)}{u}\,\dd u ,
\end{aligned}
\end{equation}
where $L$, $\fe{\cdot}{\cdot}$ and $\mathcal E^\pm$ are those of branch $\pm$, $\vth=s$ in units with $v_A=1$, and the second form of $\Hph^\pm$ uses (\ref{inv:pv}) below. The normalisations are CMM's: $W^\pm$ carries the factor $n_iT_i$ and $\Gamma^\pm$ does not, so that $n_iT_i\Gamma^\pm$ has the dimensions of $\vth W^\pm$. CMM's (B14) and (B19) for the Laguerre sectors with $l\ge2$ are the members $\varphi=1$ and $1/u$ at $\Lam=\infty$, summed over $l$. The Hermite forms $\fe{G}{G}=\sum_mp_mg_m^2$ and $\fe{G}{LG}=g^{\mathsf T}\Pw Ag$ give CMM's (B31) and (B32), with the prefactors of (\ref{eq:WGH}); on $M+1$ retained moments with odd $M$, the Hermite matrix of the helicity is the inverse of the ladder matrix at $\Lam=\infty$ (section~\ref{sec:truncated}).

The kernel of $\Hph^\pm$ contains no $\Lam$, although $L$ does. If $G/u$ has finite free energy, this follows from (\ref{eq:LUP}): then $PL^{-1}=U^{-1}$, so that $\fe{G}{L^{-1}G}=\int G^2(uF_M)^{-1}\,\dd u$. For a general $G$ the weight $1/u$ is a principal value. Appendix~\ref{appB:stone} shows that, for $\mathrm{Im}\,z>0$,
\begin{equation}\label{inv:stieltjes}
  \int\frac{a_u^2}{z-u}\,\dd u=\int\frac{G^2}{(z-u)\,F_M}\,\dd u-\frac{z}{\Lam}\,\frac{C_G(z)^2}{\diel(z)},\qquad C_G(z)=\int\frac{G(v)}{z-v}\,\dd v .
\end{equation}
At $z=\ii\eta$ with $\eta\to0^+$ the last term vanishes when $G$ is smooth in $u$ near $u=0$, as in section~\ref{sl:special}, because $C_G(\ii\eta)$ stays finite and $\diel(\ii\eta)\to1-1/\Lam\neq0$, and the Sokhotski--Plemelj formula, applied to the two remaining integrals, gives at every point $\boldsymbol r$
\begin{equation}\label{inv:pv}
  \PVint\frac{a_u^2}{u}\,\dd u=\PVint\frac{G^2}{uF_M}\,\dd u,\qquad a_0^2=\frac{G(0)^2}{F_M(0)} .
\end{equation}

CMM's gyrokinetic helicity (their 2.15) has the weight $1/(v_\parallel-\ii0)$, and the imaginary part of their Plemelj split (2.17) comes from $v_\parallel=0$. The branch analogue is the member $\varphi(u)=1/(u-\ii0)=\mathrm P(1/u)+\ii\pi\delta(u)$ of the family, and by (\ref{eq:aunit}) and (\ref{inv:pv})
\begin{equation}\label{eq:Hcomplex}
  \Hph^\pm_{\mathrm c}\equiv\frac{n_iT_i}{2\vth}\Big\langle\int\frac{(G^\pm)^2}{(u-\ii0)F_M}\,\dd u\Big\rangle
  =\Hph^\pm+\ii\,\frac{\pi n_iT_i}{2\vth}\,\Big\langle\frac{G^\pm(\boldsymbol r,0,t)^2}{F_M(0)}\Big\rangle .
\end{equation}
The imaginary part, $\pi n_iT_i\,\mathcal E^\pm(0)/(2\vth)$, is proportional to the energy of the slice at $u=0$, which is advected by $\uE$ alone (section~\ref{sec:slices}) and conserved by (\ref{inv:energy}). This is the KRMHD form of CMM's conservation of $\mathrm{Im}\,\Hph$ at each $v_\perp$ (their 2.23).

\subsection{Non-quadratic and cross-sector invariants}\label{inv:casimirs}

Each slice is a Lagrangian invariant of the volume-preserving flow of $\boldsymbol U_u$, so for every $u$ and every function $\mathcal F$ for which the integral exists
\begin{equation}\label{eq:casimir}
  \frac{\dd}{\dd t}\int\dd^3r\,\mathcal F(a_u)=0 .
\end{equation}
By (\ref{eq:sameflow}), the slices $a^\pm_u$ of the two branches and $a^{(l)}_u=G_l(u)/F_M(u)^{1/2}$ of the Laguerre sectors with $l\ge2$ share this flow, so that
\begin{equation}\label{eq:crosssector}
  \frac{\dd}{\dd t}\int\dd^3r\,\mathcal F\big(a^+_u,a^-_u,a^{(2)}_u,a^{(3)}_u,\dots\big)=0,\qquad\text{for example}\qquad\frac{\dd}{\dd t}\int\dd^3r\,a^+_u\,a^-_u=0,
\end{equation}
and integrating the last relation over $u$ with any weight gives quadratic invariants that couple the two branches. Under Fourier--Galerkin truncation in $\boldsymbol r$ only the quadratic members survive (section~\ref{sec:numerics}). Two-dimensional precedents are listed in section~\ref{sl:remlit}. In three dimensions, \citet{Tassi2022} found in Hermite truncations of drift-kinetic models with an odd number of moments one field that is advected by the $\boldsymbol E\times\boldsymbol B$ flow alone \citep[see also][\S3.4]{Tassi2023} and has an infinite family of Casimirs, the truncated counterpart of the slice at $u=0$; invariants of the form (\ref{eq:casimir}) follow, unstated, from the transformed equation of \citet{Heninger2018}, whose components stream along the uniform guide field. For $u\neq0$, where the slices stream along the perturbed field lines, the non-quadratic members of (\ref{eq:casimir}) and the cross-sector invariants (\ref{eq:crosssector}) that involve the branches $G^\pm$ are new; for the sectors with $l\ge2$ alone they follow, unstated, from S09's eq.~(212).

\subsection{The local conservation law and \texorpdfstring{$\Gamma$}{Gamma} as a flux}\label{inv:flux}

The free energy is conserved locally as well, with a flux that has a simple form. Let $w=\tfrac12\fe{G}{G}$ be the free-energy density per unit $n_iT_i$, and $\gamma=\tfrac{s}{2}\fe{G}{LG}$. The product $\fe{\cdot}{\cdot}$ acts on $u$ at fixed $\boldsymbol r$, with coefficients that do not depend on $\boldsymbol r$, so the derivatives $\{\Phi,\cdot\}$ and $\gradpar$ obey the product rule, $\gradpar\fe{f}{h}=\fe{\gradpar f}{h}+\fe{f}{\gradpar h}$. Both commute with $L$, and $L$ is symmetric in $\fe{\cdot}{\cdot}$ at each point, so that $\gradpar\fe{G}{LG}=2\,\fe{G}{\gradpar LG}$. Taking the product of (\ref{eq:branch}) with $G$ therefore gives $\partial_tw+\{\Phi,w\}+\gradpar\gamma=0$. Since $\nabla\cdot\uE=\nabla\cdot\bhat=0$, $\{\Phi,w\}=\nabla\cdot(w\,\uE)$ and $\gradpar\gamma=\nabla\cdot(\gamma\,\bhat)$, and
\begin{equation}\label{eq:localflux}
  \partial_tw+\nabla\cdot\big(w\,\uE+\gamma\,\bhat\big)=0 .
\end{equation}
In slice variables, by (\ref{sl:parseval}), $w=\tfrac12\int a_u^2\,\dd u$ and $\gamma=\tfrac{s}{2}\int u\,a_u^2\,\dd u$, so that $w\,\uE+\gamma\,\bhat=\tfrac12\int a_u^2\,\boldsymbol U_u\,\dd u$.

Each slice thus carries its free-energy density $\tfrac12a_u^2$ at its own velocity $\boldsymbol U_u$. Since $\uE$ has no $z$ component and $\hat b_z=1$, the $z$ component of the volume-averaged flux is $\langle\gamma\rangle$, which by (\ref{eq:WGH}) is $\Gamma^\pm$:
\begin{equation}\label{eq:Gammaflux}
  \Gamma^\pm=\zhat\cdot\big\langle w\,\uE+\gamma\,\bhat\big\rangle,\qquad
  \frac{n_iT_i\,\Gamma^\pm}{W^\pm}=\vth\,\frac{\int u\,\mathcal E^\pm(u)\,\dd u}{\int\mathcal E^\pm(u)\,\dd u}.
\end{equation}
Thus $n_iT_i\Gamma^\pm$ is the net flux of the free energy $W^\pm$ along $\boldsymbol B_0$, that is, the free-energy flux per unit area through the planes $z=\mathrm{const}$, averaged over $z$. The perpendicular components $\langle w\,\uE+\gamma\,\boldsymbol b_\perp\rangle$ need not vanish, and $\Gamma^\pm$ does not determine them. The ratio $n_iT_i\Gamma^\pm/W^\pm$ is the energy-weighted mean parallel velocity of the slices. CMM derived the conservation of $\Gamma^\pm$ from the moment equations (their B20--B26); its reading as a field-aligned flux of free energy is new.
